\begin{problem}{Consecutive Harshad Numbers}

    A positive integer \(n\) is called a \textit{Harshad number} if \(n\) is divisible by the sum of its
    decimal digits. Formally, writing \(S(n)\) for the sum of the digits of \(n\), we require \(S(n) \mid n\). For example, \(18\) is a Harshad number since \(S(18) = 9\) and \(9 \mid 18\), while \(19\) is \emph{not}, since \(S(19) = 10\) and \(10 \nmid 19\).

    Find the smallest \(n\) such that \(n\), \(n+1\), \(n+2\), \(n+3\) and \(n+4\) are all Harshad numbers.

    \end{problem}
